\documentclass[aspectratio=43]{beamer}

\usepackage{multicol}
\usepackage{graphicx}
\usepackage{svg}
\usepackage{fontspec}
\usepackage{csquotes}
\usepackage[main=french]{babel}
\usepackage[backend=biber]{biblatex}
\usepackage[french]{isodate}

\addbibresource{./assets/bib.bib}

\usetheme[
    background=dark,
    title page=gotham splitvert,
    numbering=pagenumber,
    sectiontocframe default=off,
    % edging default=on,
]{CambridgeUS}


\title[CERN]{L'Organisation européenne pour la recherche nucléaire (CERN)}
\author[Perry]{Aden Perry}
\institute{R.I.T.}
\date{\today}

\begin{document}

% Discuss why I chose this
%
% Je suis très interessante en le physique, et especiallemen le physique
% nucléaire. J'ai pensé qu'il serait amusant appendre plus de physique avec
% recherche de CERN.
\maketitle



% Summary
\begin{frame}{Aperçu}
	\begin{itemize}
		\item Quoi?
		      % C'est un organisation qui a l'objet de financer et améliorer
		      % l'infrastructure de physique particule de énergie haut.

		\item Où? \footfullcite{britannica}
		      % Il est situé en France et en Suisse. Le siège est près Geneva,
		      % mais il est seulment 100 hectare en Suisse. Il est plus que 450
		      % hectares en France.

		\item Quand?
		      % Il a été proposé en 1949 par Louis de Broglie à un conférence.
		      % À cause de cette proposition un conseil a été créé en 1952. 2
		      % ans plus tard, en 1954, le nom a changé à Organisation, et le
		      % version moderne de CERN etait formé.

		\item Pourquoi? \footfullcite{cern-org}
		      % CERN a été créé pour améliorer l'état de science européenne.
		      % Après le deuxième guerre mondiale, Europe n'avait pas les 
		      % meilleurs scientifiques. Beaucoup des scientifiques a migré aux
		      % États-Unis pendant la guerre.
	\end{itemize}
\end{frame}



% What is CERN known for?
\begin{frame}{Réussites}
	\begin{multicols}{2}
		\begin{itemize}
			\item Antihydrogène \footfullcite{nature}
			      % Antihydrogène est la version antimatère de hydrogène. CERN
			      % a créé antihydrogène et le conservait pour plus que 15
			      % minutes.

            \item Boson de Higgs \cite{20121}
			      % CERN a découvert un particule qui ressemble à le boson de
                  % Higgs. L'existence de le boson de Higgs confirmer
                  % l'existence du champ de Higgs, qui est important dans la
                  % modèle standard de la physique des particules.

            \item Grand collisionneur de hadrons (LHC) \footfullcite{cern-lhc}
			      % Le Grand collisionneur de hadrons est l'accélérateur des
                  % particules le plus grand et le plus puissant dans la monde.
                  % Il était utilisé, par example, pour découver le boson de
                  % Higgs.

			\item Internet \footfullcite{cern-www}
			      % En 1989, un scientifique qui travaillait à CERN a développé
                  % un logiciel qui pourrait distruber l'informatique
                  % automatiquement. Cet logiciel devenait l'Internet.
		\end{itemize}

		\columnbreak

		\center
		\includesvg[width=\linewidth]{assets/lhc.svg}
	\end{multicols}
\end{frame}



% CERN Members
\begin{frame}{Current Member States}
	% Include picture?
	\begin{itemize}
		\item 25 total member states.
		\item France contributed 159 725 200 CHF for 2026, 12.96\% of CERN's 2026 revenue. Comfortably in third, beat out by UK and Germany. Next best is 7\%, by Spain. \footfullcite{cern-finance}
	\end{itemize}
\end{frame}




% French contributions
\begin{frame}{French contributions}
\end{frame}



% Credit for presentation style
\begin{frame}{Note}
	Les manuels \textsc{Gotham} \footfullcite{gotham} et \texttt{biblatex}
	\footfullcite{biblatex} étaient consulté pour créer cette présentation.
\end{frame}



% References
\section{Références}
\begin{frame}[allowframebreaks]{Références}
	\printbibliography[heading=none]
\end{frame}

\end{document}
